\documentclass[11pt,a4paper,UTF8]{ctexart}
\usepackage{geometry}
\geometry{left=18mm,right=18mm,top=24mm,bottom=24mm,headheight=22pt,footskip=12mm}
\usepackage{amsmath,amssymb,mathtools,bm,esint,extarrows}
\usepackage{xcolor}
\usepackage{tcolorbox}
\tcbuselibrary{skins,breakable}
\usepackage{fancyhdr}
\usepackage{titlesec}
\usepackage{hyperref}
\usepackage{tikz}
\usepackage{booktabs}
\usepackage{tabularx}

% --- Color Palette (Purple & DeepPink Academic Theme) ---
\definecolor{DeepPurple}{HTML}{4C1D95}
\definecolor{TitlePurple}{HTML}{581C87}
\definecolor{SubPurple}{HTML}{6B21A8}
\definecolor{DeepPink}{HTML}{FF1493}
\definecolor{LightPinkBg}{HTML}{FFF5F8}
\definecolor{LightPurpleBg}{HTML}{F8F5FF}
\definecolor{GoldAccent}{HTML}{D97706}
\definecolor{DarkText}{HTML}{1F2937}

\hypersetup{
    colorlinks=true,
    linkcolor=TitlePurple,
    citecolor=DeepPink,
    urlcolor=SubPurple,
    pdfborder={0 0 0}
}

% --- Typography & Spacing ---
\linespread{1.3}
\setlength{\parskip}{0.35em plus 0.1em minus 0.1em}
\setlength{\parindent}{0pt}

% --- Section Titles Styling ---
\titleformat{\section}
  {\Large\bfseries\color{TitlePurple}}{\thesection}{1em}{}
  [\vspace{1mm}{\color{TitlePurple!30}\hrule height 1pt}]
\titleformat{\subsection}
  {\large\bfseries\color{SubPurple}}{\thesubsection}{0.8em}{}
\titleformat{\subsubsection}
  {\normalsize\bfseries\color{DeepPurple}}{\thesubsubsection}{0.6em}{}

% --- tcolorbox Environments ---
% 1. Theorem / Formula / Method Box (Deep Purple)
\newtcolorbox{thmbox}[1][]{
    enhanced,
    breakable,
    colback=LightPurpleBg,
    colframe=TitlePurple,
    arc=3pt,
    boxrule=0pt,
    leftrule=3.5pt,
    left=10pt,
    right=10pt,
    top=8pt,
    bottom=8pt,
    title=#1,
    coltitle=white,
    colbacktitle=TitlePurple,
    fonttitle=\bfseries\small,
    attach boxed title to top left={yshift=-2mm, xshift=4mm},
    boxed title style={boxrule=0pt, arc=2pt}
}

% 2. Solution / Formula / Derivation Box (DeepPink)
\newtcolorbox{solbox}[1][]{
    enhanced,
    breakable,
    colback=LightPinkBg,
    colframe=DeepPink,
    arc=3pt,
    boxrule=0pt,
    leftrule=3.5pt,
    left=10pt,
    right=10pt,
    top=8pt,
    bottom=8pt,
    title=#1,
    coltitle=white,
    colbacktitle=DeepPink,
    fonttitle=\bfseries\small,
    attach boxed title to top left={yshift=-2mm, xshift=4mm},
    boxed title style={boxrule=0pt, arc=2pt}
}

% 3. Learning Goals / Summary Box
\newtcolorbox{targetbox}[1][]{
    enhanced,
    breakable,
    colback=LightPurpleBg,
    colframe=DeepPurple,
    arc=3pt,
    boxrule=1pt,
    left=10pt,
    right=10pt,
    top=8pt,
    bottom=8pt,
    title=#1,
    coltitle=white,
    colbacktitle=DeepPurple,
    fonttitle=\bfseries\small,
    attach boxed title to top left={yshift=-2mm, xshift=4mm},
    boxed title style={boxrule=0pt, arc=2pt}
}

\begin{document}

% ------------------- 讲义封面 (Cover Page) -------------------
\begin{titlepage}
\thispagestyle{empty}
\begin{tikzpicture}[remember picture,overlay]
    \fill[DeepPurple] (current page.north west) rectangle ([yshift=-7cm]current page.north east);
    \draw[DeepPink, line width=3.5pt] ([yshift=-7cm]current page.north west) -- ([yshift=-7cm]current page.north east);
    \draw[GoldAccent, line width=1pt] ([yshift=-7.12cm]current page.north west) -- ([yshift=-7.12cm]current page.north east);
    
    \node[anchor=north east, opacity=0.12, text=white] at ([xshift=-1.5cm, yshift=-0.5cm]current page.north east) {
        \fontsize{70}{70}\selectfont\bfseries 2027
    };
\end{tikzpicture}

\vspace*{0.8cm}
\begin{center}
    {\color{white}\fontsize{26}{32}\selectfont\bfseries 2027 考研数学(二) 核心讲义}\\[0.6cm]
    {\color{white}\fontsize{13}{18}\selectfont\bfseries 全国硕士研究生招生考试 · 统考数学(二) 核心精讲全书}\\[1.5cm]
\end{center}

\vspace{3.5cm}
\begin{center}
    {\color{GoldAccent}\large\bfseries $\blacktriangleright$ 基础强化 · 核心题解 · 考点深水区 $\blacktriangleleft$}\\[0.8cm]
    {\color{TitlePurple}\fontsize{24}{28}\selectfont\bfseries 一元积分学的应用（积分等式和积分不等式）· 强化与专题扩展}\\[0.6cm]
    {\color{DeepPink}\large\bfseries —— 考研高阶冲刺与深水区专攻 ——}\\[1.5cm]
\end{center}

\vfill

\begin{center}
\begin{tcolorbox}[
    colback=white,
    colframe=DeepPurple,
    width=0.88\textwidth,
    arc=4pt,
    boxrule=1.2pt,
    halign=center,
    drop shadow=gray!30
]
    \vspace{3mm}
    {\bfseries\large 编著团队：EscoffierZhou}\\[2.5mm]
    {\bfseries\color{TitlePurple} 目标院校：中国科学院大学 (UCAS) 计算机专硕 (22408)}\\[2.5mm]
    {\small\color{gray} 备考铁律：手算真题数字 · 错题白纸重做 · 408 客观题为王}\\[2.5mm]
    {\small 编制版本：2027 届首发版 · \today}
    \vspace{2mm}
\end{tcolorbox}
\end{center}
\vspace{0.8cm}
\end{titlepage}

% ------------------- 目录与页面主体配置 -------------------
\pagestyle{fancy}
\fancyhf{}
\fancyhead[L]{\small\color{gray}2027 考研数学(二) · 核心讲义系列}
\fancyhead[R]{\small\color{TitlePurple}\nouppercase{\leftmark}}
\fancyfoot[C]{\small\color{gray}— 第 \thepage\ 页 —}
\fancyfoot[R]{\small\color{gray}UCAS 22408}
\renewcommand{\headrulewidth}{0.5pt}

\pagenumbering{Roman}
\tableofcontents
\newpage
\pagenumbering{arabic}


---


\section{模块一：柯西-施瓦茨（Cauchy-Schwarz）积分不等式及其深度推论}


\subsection{1. 核心定理与代数二次型判别式推导}
\begin{solbox}[分步推导与答题规范]
\textbf{[1]}  \textbf{定理内容}：
设实值函数 $f(x), g(x)$ 在区间 $[a, b]$ 上平方可积（或连续），则恒有：

\begin{equation*}
\left( \int\_a^b f(x)g(x)\,\mathrm{d}x \right)^2 \le \left( \int\_a^b f^2(x)\,\mathrm{d}x \right) \left( \int\_a^b g^2(x)\,\mathrm{d}x \right)
\end{equation*}

等号成立当且仅当存在不全为零的实数 $\lambda, \mu$ 使得 $\lambda f(x) + \mu g(x) = 0$ 几乎处处成立（若连续则恒成立，即两函数线性相关）。

\textbf{[2]}  \textbf{判别式法经典证明}：
引入实变量 $t \in \mathbb{R}$，考虑显然非负的被积函数 $[t f(x) + g(x)]^2 \ge 0$。由定积分保号性：

\begin{equation*}
P(t) = \int\_a^b [t f(x) + g(x)]^2\,\mathrm{d}x \ge 0 \quad (\forall t \in \mathbb{R})
\end{equation*}

将积分在方括号内展开并整理为关于 $t$ 的实二次三项式：

\begin{equation*}
P(t) = A t^2 + 2 B t + C \ge 0
\end{equation*}

其中：

\begin{equation*}
A = \int\_a^b f^2(x)\,\mathrm{d}x,\quad B = \int\_a^b f(x)g(x)\,\mathrm{d}x,\quad C = \int\_a^b g^2(x)\,\mathrm{d}x
\end{equation*}

若 $A = 0$，由连续函数平方积分为零知 $f(x) \equiv 0$，此时不等式两端均为 $0$，结论自明；\\ \\
若 $A > 0$，二次函数 $P(t)$ 在实数轴上恒非负，其图像开口向上且至多与横轴有一个交点。因此判别式必满足：

\begin{equation*}
\Delta = (2B)^2 - 4AC \le 0 \implies B^2 \le AC
\end{equation*}

代入系数定义即得证。
\end{solbox}


\subsection{2. 考研经典三大变式与构造技巧}
\begin{solbox}[分步推导与答题规范]
\textbf{[1]}  \textbf{倒数积型不等式（调和平均下界）}：
设 $f(x) > 0$ 且连续，令 $g(x) = \frac{1}{\sqrt{f(x)}}$，代入柯西不等式：

\begin{equation*}
\left( \int\_a^b 1\,\mathrm{d}x \right)^2 = \left( \int\_a^b \sqrt{f(x)} \cdot \frac{1}{\sqrt{f(x)}}\,\mathrm{d}x \right)^2 \le \left( \int\_a^b f(x)\,\mathrm{d}x \right) \left( \int\_a^b \frac{1}{f(x)}\,\mathrm{d}x \right)
\end{equation*}

即得考研高频结论：

\begin{equation*}
\left( \int\_a^b f(x)\,\mathrm{d}x \right) \left( \int\_a^b \frac{1}{f(x)}\,\mathrm{d}x \right) \ge (b - a)^2
\end{equation*}


\textbf{[2]}  \textbf{导数平方约束下的函数值控制}：
若 $f(a) = 0$，由微积分基本定理 $f(x) = \int_a^x f'(t)\,\mathrm{d}t$。令被积项中另一函数为常数 $1$：

\begin{equation*}
f^2(x) = \left( \int\_a^x 1 \cdot f'(t)\,\mathrm{d}t \right)^2 \le \left( \int\_a^x 1^2\,\mathrm{d}t \right) \left( \int\_a^x [f'(t)]^2\,\mathrm{d}t \right) = (x - a)\int\_a^x [f'(t)]^2\,\mathrm{d}t
\end{equation*}

两端再次在 $[a, b]$ 上对 $x$ 积分：

\begin{equation*}
\int\_a^b f^2(x)\,\mathrm{d}x \le \int\_a^b (x-a)\left(\int\_a^x [f'(t)]^2\,\mathrm{d}t\right)\,\mathrm{d}x \le \left(\int\_a^b (x-a)\,\mathrm{d}x\right)\left(\int\_a^b [f'(t)]^2\,\mathrm{d}t\right) = \frac{(b-a)^2}{2}\int\_a^b [f'(x)]^2\,\mathrm{d}x
\end{equation*}


\textbf{[3]}  \textbf{加权柯西不等式}：
若权函数 $w(x) > 0$，则：

\begin{equation*}
\left( \int\_a^b f(x)g(x)w(x)\,\mathrm{d}x \right)^2 \le \left( \int\_a^b f^2(x)w(x)\,\mathrm{d}x \right) \left( \int\_a^b g^2(x)w(x)\,\mathrm{d}x \right)
\end{equation*}

\end{solbox}

---


\section{模块二：赫尔德（Hölder）不等式与闵可夫斯基（Minkowski）不等式}


\subsection{1. 赫尔德积分不等式}
\begin{solbox}[分步推导与答题规范]
\textbf{[1]}  \textbf{定理陈述}：
设 $p, q > 1$ 且满足共轭指数关系 $\frac{1}{p} + \frac{1}{q} = 1$。对可积函数 $f(x), g(x)$，恒有：

\begin{equation*}
\int\_a^b |f(x)g(x)|\,\mathrm{d}x \le \left( \int\_a^b |f(x)|^p\,\mathrm{d}x \right)^{1/p} \left( \int\_a^b |g(x)|^q\,\mathrm{d}x \right)^{1/q}
\end{equation*}

当 $p = q = 2$ 时，赫尔德不等式退化为柯西-施瓦茨不等式。

\textbf{[2]}  \textbf{基于杨氏（Young）不等式的微积分证明}：
由杨氏不等式：对任意 $u, v \ge 0$，有 $u v \le \frac{u^p}{p} + \frac{v^q}{q}$。\\ \\
引入归一化量：

\begin{equation*}
u(x) = \frac{|f(x)|}{\left( \int\_a^b |f|^p \right)^{1/p}},\quad v(x) = \frac{|g(x)|}{\left( \int\_a^b |g|^q \right)^{1/q}}
\end{equation*}

代入杨氏不等式并在 $[a, b]$ 上逐项积分：

\begin{equation*}
\int\_a^b u(x)v(x)\,\mathrm{d}x \le \frac{1}{p}\int\_a^b u^p(x)\,\mathrm{d}x + \frac{1}{q}\int\_a^b v^q(x)\,\mathrm{d}x = \frac{1}{p} \cdot 1 + \frac{1}{q} \cdot 1 = 1
\end{equation*}

整理即得结论。
\end{solbox}


\subsection{2. 闵可夫斯基三角不等式}
\begin{solbox}[分步推导与答题规范]
\textbf{[1]}  \textbf{范数三角不等式形式}：
对 $p \ge 1$，恒有：

\begin{equation*}
\left( \int\_a^b |f(x) + g(x)|^p\,\mathrm{d}x \right)^{1/p} \le \left( \int\_a^b |f(x)|^p\,\mathrm{d}x \right)^{1/p} + \left( \int\_a^b |g(x)|^p\,\mathrm{d}x \right)^{1/p}
\end{equation*}

在泛函分析与考研分析题中，该定理确立了 $L^p[a, b]$ 空间满足范数三角不等式（即 $\|f+g\|_p \le \|f\|_p + \|g\|_p$）。
\end{solbox}

---


\section{模块三：凸函数不等式体系：琴生（Jensen）与 Hermite-Hadamard 不等式}


\subsection{1. 连续型琴生不等式}
\begin{solbox}[分步推导与答题规范]
\textbf{[1]}  \textbf{理论结构}：
设 $\varphi(u)$ 是区间 $I$ 上的下凸函数（若二阶可导则 $\varphi''(u) \ge 0$）。若权函数 $p(x) \ge 0$ 满足归一化条件 $\int_a^b p(x)\,\mathrm{d}x = 1$，且 $f(x)$ 的值域落在 $I$ 内，则：

\begin{equation*}
\varphi\left( \int\_a^b f(x)p(x)\,\mathrm{d}x \right) \le \int\_a^b \varphi(f(x))p(x)\,\mathrm{d}x
\end{equation*}


\textbf{[2]}  \textbf{无偏权重特例}：
取 $p(x) = \frac{1}{b-a}$，可得均值形式：

\begin{equation*}
\varphi\left( \frac{1}{b-a}\int\_a^b f(x)\,\mathrm{d}x \right) \le \frac{1}{b-a}\int\_a^b \varphi(f(x))\,\mathrm{d}x
\end{equation*}


\textbf{[3]}  \textbf{极具实战价值的特例}：
\textbf{(1)}  取 $\varphi(u) = u^2$：$\left( \int_a^b f(x)\,\mathrm{d}x \right)^2 \le (b-a)\int_a^b f^2(x)\,\mathrm{d}x$；
\textbf{(2)}  取 $\varphi(u) = \mathrm{e}^u$：$\exp\left( \frac{1}{b-a}\int_a^b f(x)\,\mathrm{d}x \right) \le \frac{1}{b-a}\int_a^b \mathrm{e}^{f(x)}\,\mathrm{d}x$（连续型算术-几何平均值不等式）；
\textbf{(3)}  取 $\varphi(u) = -\ln u$（对 $u > 0$）：$\int_a^b \ln f(x)\,\mathrm{d}x \le (b-a)\ln\left(\frac{1}{b-a}\int_a^b f(x)\,\mathrm{d}x\right)$。
\end{solbox}


\subsection{2. Hermite-Hadamard 双向夹逼不等式}
\begin{solbox}[分步推导与答题规范]
\textbf{[1]}  \textbf{定理表述}：
若 $f(x)$ 在 $[a, b]$ 上是下凸函数，则中点函数值、区间平均积分值与端点算术平均值构成三级偏序：

\begin{equation*}
f\left( \frac{a+b}{2} \right) \le \frac{1}{b-a}\int\_a^b f(x)\,\mathrm{d}x \le \frac{f(a) + f(b)}{2}
\end{equation*}


\textbf{[2]}  \textbf{几何推导直观}：
左半侧由切线支撑性质保证（切线不等式积分，过中点的切线在曲线下方）；
右半侧由割线凸性保证（连接 $(a, f(a))$ 与 $(b, f(b))$ 的割线在曲线上方，梯形面积大于曲线下面积）。
\end{solbox}

---


\section{模块四：维尔廷格（Wirtinger）与庞加莱（Poincaré）型能量控制不等式}


\subsection{1. 庞加莱（Poincaré）一阶能量不等式}
\begin{solbox}[分步推导与答题规范]
\textbf{[1]}  \textbf{端点为零情形}：
若 $f(a) = f(b) = 0$ 且 $f'(x)$ 连续，则存在依赖于区间长度的常数 $C$，使得：

\begin{equation*}
\int\_a^b f^2(x)\,\mathrm{d}x \le \frac{(b-a)^2}{\pi^2}\int\_a^b [f'(x)]^2\,\mathrm{d}x
\end{equation*}


\textbf{[2]}  \textit{}初等微分法证明弱形式（系数为 $\frac{(b-a)^2}{2}$）\textit{}：
由前文柯西推论，对任意 $x \in [a, b]$，有 $f^2(x) \le (x-a)\int_a^b [f'(t)]^2\,\mathrm{d}t$。两边在 $[a, b]$ 上积分即可直接导出弱形式：

\begin{equation*}
\int\_a^b f^2(x)\,\mathrm{d}x \le \frac{(b-a)^2}{2}\int\_a^b [f'(x)]^2\,\mathrm{d}x
\end{equation*}

\end{solbox}


\subsection{2. 维尔廷格（Wirtinger）零均值不等式}
\begin{solbox}[分步推导与答题规范]
\textbf{[1]}  \textbf{零均值条件约束}：
若 $f(x)$ 满足 $\int_a^b f(x)\,\mathrm{d}x = 0$，且 $f'(x) \in L^2[a, b]$，则恒有：

\begin{equation*}
\int\_a^b f^2(x)\,\mathrm{d}x \le \left( \frac{b-a}{2\pi} \right)^2 \int\_a^b [f'(x)]^2\,\mathrm{d}x
\end{equation*}

该不等式在物理中表征驻波与振动的最小动能约束，在考研中常作为压轴分析题的背景原型。
\end{solbox}

---


\section{模块五：变上限积分方程向常微分方程的转化技术}


\subsection{1. 积分方程的分类与处理哲学}
\begin{solbox}[分步推导与答题规范]
\textbf{[1]}  \textbf{微分降阶法（针对变上限微分-积分方程）}：
当方程中出现自变量 $x$ 既在被积函数内部又在积分上限中时，必须先将内外部变量剥离解耦，再逐次求导转化为常微分方程初值问题。
\textbf{[2]}  \textbf{待定常数法（针对 Fredholm 型常上下限积分方程）}：
当方程形如 $f(x) = g(x) + \lambda \int_a^b K(x, t)f(t)\,\mathrm{d}t$ 时，若核可分离 $K(x, t) = \sum u_i(x)v_i(t)$，由于定积分值是纯常数，直接设 $\int_a^b v_i(t)f(t)\,\mathrm{d}t = C_i$ 转化为代数方程组求解。
\end{solbox}


\subsection{2. 经典模型解题范式全剖析}
\begin{solbox}[分步推导与答题规范]
\textbf{[1]}  \textbf{卷积型变上限积分方程转化}：
求解方程：

\begin{equation*}
f(x) = x + \int\_0^x (x - t)f(t)\,\mathrm{d}t
\end{equation*}

解法步骤：
第一步（拆分积分并求一阶导）：

\begin{equation*}
f(x) = x + x \int\_0^x f(t)\,\mathrm{d}t - \int\_0^x t f(t)\,\mathrm{d}t
\end{equation*}

两端对 $x$ 求导：

\begin{equation*}
f'(x) = 1 + \left( \int\_0^x f(t)\,\mathrm{d}t + x f(x) \right) - x f(x) = 1 + \int\_0^x f(t)\,\mathrm{d}t
\end{equation*}

由原方程知 $f(0) = 0$；由导数式知 $f'(0) = 1$。\\ \\
第二步（求二阶导转化为二阶常系数齐次线性微分方程）：

\begin{equation*}
f''(x) = f(x) \implies f''(x) - f(x) = 0
\end{equation*}

第三步（解微分方程定特解）：
特征方程为 $r^2 - 1 = 0 \implies r = \pm 1$。通解为：

\begin{equation*}
f(x) = C\_1 \mathrm{e}^x + C\_2 \mathrm{e}^{-x}
\end{equation*}

代入初值条件 $f(0) = 0 \implies C_1 + C_2 = 0$；$f'(0) = 1 \implies C_1 - C_2 = 1$。\\ \\
解得 $C_1 = \frac{1}{2}, C_2 = -\frac{1}{2}$。\\ \\
因此：

\begin{equation*}
f(x) = \frac{\mathrm{e}^x - \mathrm{e}^{-x}}{2} = \sinh x
\end{equation*}

\end{solbox}

---


\section{模块六：区间再现与对称变换专题技巧矩阵}


\subsection{1. 经典三大对称积分换元}
\begin{solbox}[分步推导与答题规范]
\textbf{[1]}  \textbf{闭区间对称换元（倒角公式）}：
对任意区间 $[a, b]$：

\begin{equation*}
\int\_a^b f(x)\,\mathrm{d}x = \int\_a^b f(a + b - x)\,\mathrm{d}x = \frac{1}{2}\int\_a^b [f(x) + f(a + b - x)]\,\mathrm{d}x
\end{equation*}

若被积函数满足 $f(a + b - x) + f(x) = C$（关于点 $\left(\frac{a+b}{2}, \frac{C}{2}\right)$ 中心对称），则定积分直接秒杀：

\begin{equation*}
\int\_a^b f(x)\,\mathrm{d}x = \frac{C(b - a)}{2}
\end{equation*}


\textbf{[2]}  \textbf{三角周期对称换元}：
在 $[0, \pi]$ 上令 $x = \pi - t$：

\begin{equation*}
\int\_0^\pi x g(\sin x)\,\mathrm{d}x = \frac{\pi}{2}\int\_0^\pi g(\sin x)\,\mathrm{d}x
\end{equation*}

在 $[0, \pi/2]$ 上令 $x = \pi/2 - t$：

\begin{equation*}
\int\_0^{\pi/2} \frac{\sin^n x}{\sin^n x + \cos^n x}\,\mathrm{d}x = \frac{\pi}{4}
\end{equation*}


\textbf{[3]}  \textbf{倒数对称换元（乘积反演）}：
令 $x = \frac{1}{t}$，可将 $[0, 1]$ 与 $[1, +\infty]$ 建立对偶联系，用于化简具有 $x \leftrightarrow \frac{1}{x}$ 反演不变性的积分核。
\end{solbox}

---


\section{模块七：格隆沃尔（Grönwall）积分不等式与渐近估计}


\subsection{1. 定理陈述与微分因子推导}
\begin{solbox}[分步推导与答题规范]
\textbf{[1]}  \textbf{微分方程估计的核心工具}：
设连续函数 $u(t) \ge 0$，且对常数 $C \ge 0$ 和非负连续函数 $k(t)$，满足积分不等式：

\begin{equation*}
u(t) \le C + \int\_0^t k(s)u(s)\,\mathrm{d}s \quad (\forall t \ge 0)
\end{equation*}

则 $u(t)$ 必满足指数上界估计：

\begin{equation*}
u(t) \le C \exp\left( \int\_0^t k(s)\,\mathrm{d}s \right)
\end{equation*}

特别地，若 $C = 0$，则必有 $u(t) \equiv 0$。

\textbf{[2]}  \textbf{构造一阶微分不等式证明法}：
令变上限函数 $v(t) = C + \int_0^t k(s)u(s)\,\mathrm{d}s$，则 $v(0) = C$ 且 $u(t) \le v(t)$。\\ \\
对 $v(t)$ 求导：

\begin{equation*}
v'(t) = k(t)u(t) \le k(t)v(t) \implies v'(t) - k(t)v(t) \le 0
\end{equation*}

两边同乘积分因子 $\exp\left(-\int_0^t k(s)\,\mathrm{d}s\right) > 0$：

\begin{equation*}
\frac{\mathrm{d}}{\mathrm{d}t}\left[ v(t)\exp\left(-\int\_0^t k(s)\,\mathrm{d}s\right) \right] \le 0
\end{equation*}

从 $0$ 到 $t$ 积分：

\begin{equation*}
v(t)\exp\left(-\int\_0^t k(s)\,\mathrm{d}s\right) - v(0) \le 0 \implies v(t) \le C \exp\left(\int\_0^t k(s)\,\mathrm{d}s\right)
\end{equation*}

由于 $u(t) \le v(t)$，不等式得证。
\end{solbox}


\subsection{2. 考研经典应用：微分方程解的唯一性证明}
\begin{solbox}[分步推导与答题规范]
\textbf{[1]}  \textbf{范式应用}：
对于初值问题 $y' = f(x, y), y(x_0) = y_0$，若 $f$ 满足李普希茨条件 $|f(x, y_1) - f(x, y_2)| \le L |y_1 - y_2|$。\\ \\
若存在两个解 $y_1(x), y_2(x)$，定义偏差 $u(x) = |y_1(x) - y_2(x)|$。则：

\begin{equation*}
u(x) = \left| \int\_{x\_0}^x [f(t, y\_1(t)) - f(t, y\_2(t))]\,\mathrm{d}t \right| \le L \int\_{x\_0}^x u(t)\,\mathrm{d}t
\end{equation*}

由格隆沃尔不等式（取 $C = 0, k(t) = L$），立得 $u(x) \equiv 0 \implies y_1(x) \equiv y_2(x)$，解必唯一。
\end{solbox}

\end{document}
